% !TeX root = ../../progress-sem1.tex
% 2.x สรุปการศึกษาค้นคว้า
\section{สรุปการศึกษาค้นคว้า}

% \begin{table}[H]
%   \centering
%   \small
%   \caption{สรุปประเด็นจากงานวิจัยและเครื่องมือที่เกี่ยวข้อง}
%   \label{tab:literature_summary}
%   \begin{tabularx}{\textwidth}{
%       >{\raggedright\arraybackslash}p{0.18\textwidth}
%       >{\raggedright\arraybackslash}p{0.22\textwidth}
%       >{\raggedright\arraybackslash}X
%       >{\raggedright\arraybackslash}X
%     }
%     \hline
%     \textbf{กลุ่มประเด็น} &
%     \textbf{งานที่เกี่ยวข้อง} &
%     \textbf{สิ่งที่มีการศึกษาแล้ว} &
%     \textbf{ช่องว่างที่พบ} \\
%     \hline

%     การสร้าง UI และต้นแบบด้วย Generative AI &
%     \cite{suh2025storyensemble, feng2023designing, yuan2024towards, kolthoff2024} &
%     Generative AI สามารถช่วยสร้าง wireframe, prototype และทางเลือกของ UI
%     จากความต้องการของนักออกแบบ เพื่อลดเวลาและแรงงานในการสร้างต้นแบบ
%     และช่วยให้สามารถสำรวจทางเลือกของการออกแบบได้มากขึ้น &
%     ผลลัพธ์ที่สร้างโดย Generative AI ยังพบปัญหาด้าน layout, style และการตีความ
%     function ที่ไม่ตรงกับความต้องการ รวมถึงยังไม่ได้เชื่อมกับกระบวนการ
%     ทดสอบและปรับปรุงจากข้อเสนอแนะของผู้ใช้อย่างครบวงจร \\

%     \hline

%     การทดสอบและการสกัดความต้องการจากต้นแบบ &
%     \cite{nielsen2000five, virzi1992refining,
%     nielsen1993mathematical, pacheco2018, lu2026} &
%     การทดสอบกับผู้ใช้กลุ่มเล็กหลายรอบสามารถช่วยค้นพบปัญหาได้อย่างมีประสิทธิภาพ
%     และการใช้ prototype ช่วยให้นักออกแบบสามารถสกัดความต้องการและรับ
%     ข้อเสนอแนะก่อนพัฒนาระบบจริงได้ &
%     การทำการทดสอบซ้ำหลายรอบยังสร้างภาระในการเตรียมหลายรูปแบบของ UI
%     รวบรวมข้อเสนอแนะ และนำผลจากแต่ละรอบกลับไปปรับปรุงการออกแบบ \\

%     \hline

%     การตีความข้อเสนอแนะและการสร้าง User Persona &
%     \cite{nngroup_satisfaction, nngroup_qualitative,
%     shukla2025, mitchell1997, shin2024} &
%     งานวิจัยแสดงให้เห็นว่าความชอบของผู้ใช้อาจไม่ตรงกับปัญหาการใช้งานจริง
%     และข้อเสนอแนะยังแตกต่างกันตามบทบาทและความต้องการของผู้ให้ข้อมูล
%     นอกจากนี้ LLM สามารถช่วยสรุปข้อมูลเพื่อสร้าง User Persona ได้
%     แต่ยังต้องอาศัยมนุษย์ในการระบุและตรวจสอบลักษณะที่สำคัญของผู้ใช้ &
%     ยังต้องมีกระบวนการเชื่อมโยงข้อเสนอแนะกับกลุ่มผู้ใช้
%     แยกความต้องการและความชอบออกจากกัน และนำข้อมูลเหล่านี้ไปสังเคราะห์
%     เป็น Persona ที่นักออกแบบสามารถตรวจสอบและนำกลับไปใช้ในการออกแบบได้ \\

%     \hline

%     การใช้ปัญญาประดิษฐ์วิเคราะห์ข้อเสนอแนะร่วมกับมนุษย์ &
%     \cite{gao2024collabcoder, gebreegziabher2023patat,
%     song2025, amershi2019, pourasad2025} &
%     LLM สามารถช่วยจัดกลุ่ม สรุป และวิเคราะห์ข้อมูลเชิงคุณภาพได้
%     แต่ยังมีข้อจำกัดด้านอคติและอาจพลาดข้อมูลสำคัญ
%     งานวิจัยจึงสนับสนุนการให้มนุษย์เข้ามาตรวจสอบและตัดสินใจร่วมกับปัญญาประดิษฐ์ &
%     จำเป็นต้องมีกระบวนการที่ให้นักออกแบบตรวจสอบ แก้ไข
%     และยืนยันผลการวิเคราะห์ของปัญญาประดิษฐ์ก่อนที่ผลลัพธ์จะถูกนำไปใช้
%     ปรับปรุงงานออกแบบ \\

%     \hline

%     เครื่องมือที่มีอยู่ในปัจจุบัน &
%     \cite{kameleoon, otterab, lyssna, useberry,
%     figma_agent, bugherd, markerio} &
%     เครื่องมือแต่ละประเภทสนับสนุนบางส่วนของกระบวนการ เช่น
%     การทดสอบหลายรูปแบบ การทดสอบ prototype การสร้าง UI ด้วย AI
%     หรือการให้ข้อเสนอแนะบนงานออกแบบ &
%     ความสามารถเหล่านี้ยังแยกอยู่ในหลายเครื่องมือ
%     และยังไม่มีการเชื่อมตั้งแต่การสร้างทางเลือกของ UI
%     การทดสอบ การวิเคราะห์ข้อเสนอแนะ การสร้าง Persona
%     จนถึงการนำผลกลับไปใช้ในรอบการออกแบบถัดไป \\

%     \hline
%   \end{tabularx}
% \end{table}

\thaijustify{
  จากการศึกษางานวิจัยและเครื่องมือที่เกี่ยวข้อง พบว่าแต่ละแนวทางสามารถสนับสนุนบางส่วนของ UX/UI DLC ได้แล้ว แต่กระบวนการตั้งแต่
  การสร้างต้นแบบ การทดสอบกับผู้ใช้ การรวบรวมและวิเคราะห์ข้อเสนอแนะจนถึงการนำข้อมูลกลับไปปรับปรุงการออกแบบยังคงแยกออกจากกัน
}
\thaijustify{
  ดังนั้น โครงงานนี้จึงมุ่งเชื่อมกระบวนการเหล่านี้เข้าด้วยกันโดยใช้ Generative AI ช่วยสร้าง two-alternative UI จากความต้องการทางการออกแบบ แล้วนำไปทำ Comparative Usability Test และให้ผู้เข้าร่วมทดสอบสามารถให้ข้อเสนอแนะลงบน component ได้โดยตรง จากนั้นระบบจะช่วยรวบรวมและวิเคราะห์ข้อเสนอแนะเพื่อค้นหาปัญหา ความต้องการ ความชอบและลักษณะร่วมของผู้ใช้แต่ละกลุ่ม เพื่อนำไปสนับสนุนการสร้าง Persona Artifact โดยให้นักออกแบบเป็นผู้ตรวจสอบและตัดสินใจผลลัพธ์ก่อนนำข้อมูลที่ผ่านการยืนยันแล้วกลับไปใช้ในการสร้างและทดสอบรอบถัดไป
}
\thaijustify{
  แนวทางดังกล่าวทำให้ โครงงานไม่ได้มุ่งใช้ Generative AI เพื่อสร้าง two-alternative UI เพียงอย่างเดียว แต่ใช้ปัญญาประดิษฐ์เพื่อสนับสนุน UX/UI DLC ตั้งแต่การสร้างต้นแบบ การทดสอบการใช้งาน ไปจนถึงการนำความต้องการ
  ของผู้ใช้กลับมาใช้ในการตัดสินใจเชิงออกแบบ โดยยังคงให้นักออกแบบเป็นผู้ควบคุมการตัดสินใจในแต่ละรอบ
}

% \thaijustify{
%   และจากการทบทวนวรรณกรรมทำให้เห็นช่องว่างทางการวิจัย คือ ปัจจุบันยังไม่มีเครื่องมือหรืองานวิจัยใดที่รวบรวมความสามารถทั้ง 4 ด้านไว้ในระบบเดียว ได้แก่ (1) การใช้ปัญญาประดิษฐ์สร้าง two-alternative UI ภายใต้ Design system ที่กำหนด (2) การให้ผู้ทดสอบให้ข้อเสนอแนะ (3) การใช้ปัญญาประดิษฐ์แยกแยะ pain point ออกจากความชอบร่วมกับการตัดสินใจของนักออกแบบ และ (4) การสร้าง Persona Artifact จากผลลัพธ์
% }
% \thaijustify{สิ่งที่โครงงานนี้ควรให้ความสำคัญเป็น คือ การพัฒนาแพลตฟอร์มนี้ให้ครอบคลุมกระบวนการตั้งแต่รับโจทย์ สร้างต้นแบบ รวบรวมข้อเสนอแนะ และวิเคราะห์ผลลัพธ์ เพื่อให้นักออกแบบสามารถส่งมอบงานที่มีการตรวจสอบความต้องการของผู้ใช้ที่แท้จริง พร้อมด้วย design system และต้นแบบที่ลูกค้าพึงพอใจให้นักพัฒนาต่อไปได้อย่างมีประสิทธิภาพ
% }
